\newcommand{\GBMEMRate}{0.500}
\newcommand{\GBMEMWeakRate}{0.999}
\newcommand{\GBMMilRate}{0.991}
\newcommand{\GBMMilWeakRate}{1.000}
\newcommand{\GBMAnalyticWeakRate}{0.9993}
\newcommand{\GBMFitWindow}{$N=64,128,256,512,1024$}
\newcommand{\GBMExactMean}{105.1987}
\newcommand{\GBMExactVariance}{1045.0120}
\newcommand{\GBMExactMeanCI}{[105.0403, 105.3571]}
\newcommand{\NARate}{0.491}
\newcommand{\NAFitWindow}{$N=8,16,32,64$}
\newcommand{\NAExactMean}{105.08796}
\newcommand{\NAExactMeanCI}{[104.89962, 105.27629]}
\newcommand{\NAPayoff}{14.53804}
\newcommand{\NAPayoffCI}{[14.40515, 14.67093]}
\newcommand{\NAWeakInterpretation}{At $N=8,16,32$, the intervals exclude zero, so these coarse-grid payoff differences are distinguishable from sampling noise. At $N=64,128,256$, the intervals include zero: the available sample cannot reliably resolve the sign. This does not prove zero bias, and no weak order is assigned to the complete grid range.}
\newcommand{\NAReferenceInterpretation}{The average absolute difference falls from 0.05679 for 2048 versus 4096 steps to 0.04016 for 4096 versus 8192 steps. The second difference is about 0.707 of the first, consistent with half-order scaling. Relative to the coarse-to-finest differences, 0.04016 represents 3.2\% to 8.9\% on the $N=8,16,32,64$ window, but 12.6\% and 18.0\% at $N=128,256$. Reference resolution therefore matters more in relative terms on the finer coarse grids. The principal slope uses the less reference-sensitive four-grid window. These comparisons measure sensitivity to reference refinement; they do not bound the unknown error against the exact solution.}
